\heading{实验物理中的统计方法-最小二乘法}





\begin{question}
Galileo
研究运动的实验之一是小球和斜坡的实验。在离开斜坡边缘之前，小球的轨迹变成水平，如图
7.1 所示。对于不同的高度 \(h\)，测量从斜坡边缘到落地点的水平距离
\(d\)。1608 年，Galileo 测量了 5 组数据，如表 7.1 所示。假设高度 \(h\)
的误差可以忽略，水平距离 \(d\) 可以看做独立的标准差 \(\sigma=15\) punti
的高斯随机变量。（实际上我们不清楚 Galileo 如何估计测量误差，但是 1-2\%
的误差是可以接受的。）此外，我们知道如果 \(h=0\)，则水平距离 \(d\)
将为零，即，如果球从斜坡边缘出发，它将垂直落到地上。
\pyfig[0.72\linewidth]{figures/fig_07_01_galileo_schematic.png}{图 7.1：Galileo 斜坡实验的几何关系示意图。}


\begin{center}
\begin{tabular}{cc}
\toprule
\(h\) & \(d\) \\
\midrule
1000 & 1500 \\
828 & 1340 \\
800 & 1328 \\
600 & 1172 \\
300 & 800 \\
\bottomrule
\end{tabular}
\par\vspace{0.25em}{\small 表 7.1: Galileo 斜坡实验的 5 组数据。给定初始高度 \(h\)，\(d\) 为落地前的水平距离。单位为 punti，\(1\,\mathrm{punto}\simeq1\,\mathrm{mm}\)。\par}
\end{center}

\begin{enumerate}
\item
  考虑 \(h\) 和 \(d\) 的关系为如下形式
\end{enumerate}

\[
d=\alpha h
\]

以及

\[
d=\alpha h+\beta h^2.
\]

求参数 \(\alpha\) 和 \(\beta\) 的最小二乘估计量。对应于这两个假设的最小
\(\chi^2\) 和 \(P\)-值分别为多少？

\begin{enumerate}[start=2]
\item
  假设 \(d\) 和 \(h\) 的关系为如下形式
\end{enumerate}

\[
d=\alpha h^\beta.
\]

写一段程序对 \(\alpha\) 和 \(\beta\)
进行最小二乘拟合。注意这是参数的非线性函数，必须数值求解。

\begin{enumerate}[start=3]
\item
  Galileo
  认为运动是水平分量和垂直分量的叠加，其中水平运动是匀速运动，垂直速度在斜坡的最低处为零，随后随时间线性增加。证明这将导致关系式
\end{enumerate}

\[
d=\alpha\sqrt{h}.
\]

求 \(\alpha\) 的最小二乘估计量以及最小 \(\chi^2\)。该假设的
\(P\)-值是多少？
\begin{solution}
把数据记为
\[
(h_i,d_i)=(1000,1500),(828,1340),(800,1328),(600,1172),(300,800),
\]
并把各 \(d_i\) 视作标准差为 \(\sigma=15\) punti 的独立高斯变量。

\begin{enumerate}
\item 比较模型

\begin{enumerate}
\item
  线性模型

  \[
  d=\alpha h.
  \]

  加权最小二乘等价于最小化

  \[
  \chi^2(\alpha)=\sum_i\frac{(d_i-\alpha h_i)^2}{\sigma^2}.
  \]

  令导数为零，得到

  \[
  \hat\alpha=\frac{\sum_i h_i d_i}{\sum_i h_i^2}.
  \]
\item
  二次模型

  \[
  d=\alpha h+\beta h^2.
  \]

  这是线性参数模型，可直接由正规方程或矩阵最小二乘求出
  \((\hat\alpha,\hat\beta)\)。
\end{enumerate}

程序计算结果：

\begin{itemize}
\item
  \(d=\alpha h\)： \[
  \hat\alpha=1.662756,
  \qquad
  \chi^2_{\min}=661.99,
  \qquad
  n_{\rm dof}=4,
  \qquad
  P\approx5.9\times10^{-142}.
  \]
\item
  \(d=\alpha h+\beta h^2\)： \[
  \hat\alpha=2.792921,
  \qquad
  \hat\beta=-1.35055\times10^{-3},
  \] \[
  \chi^2_{\min}=64.74,
  \qquad
  n_{\rm dof}=3,
  \qquad
  P\approx5.7\times10^{-14}.
  \]
\end{itemize}

这说明如果坚持把误差取为 \(15\) punti，那么前两个模型都拟合得很差。

\item 幂律模型

设

\[
d=\alpha h^\beta.
\]

这是非线性参数拟合。程序数值最小化得到

\[
\hat\alpha=43.7606,
\qquad
\hat\beta=0.511056,
\]

\[
\chi^2_{\min}=3.756,
\qquad
n_{\rm dof}=3,
\qquad
P=0.289.
\]

这已经是可以接受的拟合。

\item Galileo 的平方根模型

由 Galileo 的运动学推导可得

\[
d=\alpha\sqrt h.
\]

最小二乘估计为

\[
\hat\alpha=
\frac{\sum_i d_i\sqrt{h_i}}{\sum_i h_i}.
\]

数值结果：

\[
\hat\alpha=47.0858,
\qquad
\chi^2_{\min}=4.208,
\qquad
n_{\rm dof}=4,
\qquad
P=0.379.
\]

因此在这组数据上，\(d\propto\sqrt h\)
的物理模型优于简单线性模型，也优于随意加入 \(h^2\)
项后的二次模型。它与幂律模型得到的 \(\beta\approx0.51\) 彼此一致。
\pyfig[0.72\linewidth]{figures/fig_07_01_galileo_fit.png}{Galileo 数据及平方根模型、一般幂律模型的拟合曲线。}\end{enumerate}
\end{solution}
\end{question}



\begin{question}
考虑对直方图的最小二乘拟合。直方图对应区间 \(i=1,\ldots,N\) 的事例数为
\(y_i\)，理论预言值为

\[
\lambda_i(\theta)=n\int_{x_i^{\min}}^{x_i^{\max}}f(x;\theta)\,dx,
\]

其中 \(f(x;\theta)\) 依赖于未知参数 \(\theta\)。假设用参数 \(\nu\)
代替总事例数 \(n\)，并且该参数在最小化

\[
\chi^2(\theta,\nu)=\sum_{i=1}^{N}\frac{(y_i-\lambda_i(\theta,\nu))^2}{\sigma_i^2}
\]

时与其它参数同时调整。

\begin{enumerate}
\item
  证明如果 \(\sigma_i^2=\lambda_i\)，则总事例数的估计量为
\end{enumerate}

\[
\hat\nu_{\rm LS}=n+\frac{\chi^2_{\min}}{2}.
\]

\begin{enumerate}[start=2]
\item
  证明如果 \(\sigma_i^2=y_i\)（修正后的最小二乘），则估计量为
\end{enumerate}

\[
\hat\nu_{\rm MLS}=n-\chi^2_{\min}.
\]
\begin{solution}
若第 \(i\) 个 bin 的理论期望写作

\[
\lambda_i(\theta,\nu)=\nu\int_{x_i^{\min}}^{x_i^{\max}} f(x;\theta)\,dx,
\]

并把总归一化参数也一并放进拟合，那么：

\begin{enumerate}
\item 当取 \(\sigma_i^2=\lambda_i\) 时

最小化

\[
\chi^2(\theta,\nu)=\sum_i\frac{(y_i-\lambda_i)^2}{\lambda_i}
\]

对 \(\nu\) 求导可得

\[
\hat\nu_{LS}=n+\frac{\chi^2_{\min}}{2}.
\]

因此普通 LS 会把总事件数估高。

\item 当取 \(\sigma_i^2=y_i\) 时

最小化

\[
\chi^2(\theta,\nu)=\sum_i\frac{(y_i-\lambda_i)^2}{y_i}
\]

则有

\[
\hat\nu_{MLS}=n-\chi^2_{\min}.
\]

所以修正最小二乘会把总事件数估低。

这两个结果本质上都来自``把分母里的方差也当成与待估参数有关''后，目标函数和真正的对数似然不再一致。
\end{enumerate}
\end{solution}
\end{question}



\begin{question}
考虑对直方图的最小二乘拟合，第 \(i\) 个区间的事例数为
\(y_i\)，\(i=1,\ldots,N\)，对应的期待值为
\(\lambda_i(\theta)\)。假设总事例数 \(n\) 可以看作常数，于是 \(y_i\)
服从多项分布。

\begin{enumerate}
\item
  求协方差矩阵
  \(V_{ij}=\operatorname{cov}[y_i,y_j]\)。为什么该矩阵的逆不存在？
\item
  考虑只使用前 \(N-1\)
  个区间进行拟合。求协方差矩阵的逆，并证明这等价于对所有 \(N\)
  个区间进行拟合但不考虑相关性。
\end{enumerate}
\begin{solution}
设总事件数固定为 \(n\)，各 bin 概率为 \(p_i\)，故
\(y=(y_1,\ldots,y_N)\sim\operatorname{Multinomial}(n;p_1,\ldots,p_N)\)。

\begin{enumerate}
\item 协方差矩阵为
\[
V_{ij}=\operatorname{cov}(y_i,y_j)
=n p_i(\delta_{ij}-p_j).
\]
即
\[
V_{ii}=np_i(1-p_i),
\qquad
V_{ij}=-np_ip_j\quad(i\ne j).
\]
该矩阵不可逆，因为总数固定给出严格约束
\[
\sum_{i=1}^N y_i=n.
\]
于是 \(\sum_i (y_i-E[y_i])=0\)，存在一个零方差方向，矩阵秩至多为
\(N-1\)。

\item 只取前 \(N-1\) 个 bin。令 \(p_N=1-\sum_{i=1}^{N-1}p_i\)，则前
\(N-1\) 个 bin 的协方差矩阵为
\[
V=n(D-pp^T),
\qquad
D=\operatorname{diag}(p_1,\ldots,p_{N-1}).
\]
用 Sherman--Morrison 公式得
\[
(V^{-1})_{ij}=\frac1n\left(\frac{\delta_{ij}}{p_i}+\frac1{p_N}\right),
\qquad i,j=1,\ldots,N-1.
\]
设残差 \(r_i=y_i-np_i\)。由于总数固定，\(r_N=-\sum_{i=1}^{N-1}r_i\)。
于是相关协方差给出的二次型为
\[
\sum_{i,j=1}^{N-1}r_i(V^{-1})_{ij}r_j
=\sum_{i=1}^{N-1}\frac{r_i^2}{np_i}
+\frac{(\sum_{i=1}^{N-1}r_i)^2}{np_N}
=\sum_{i=1}^N\frac{r_i^2}{np_i}.
\]
这正是使用全部 \(N\) 个 bin、但只取对角方差
\(\sigma_i^2=np_i\) 时的 Pearson \(\chi^2\)。因此二者等价。
\end{enumerate}
\end{solution}
\end{question}



\begin{question}
假设样本容量为 \(n\)，利用该样本得到 \(N\)
个量的测量值：\(y_1,\ldots,y_N\)。这 \(N\)
个测量值将被用于最小二乘拟合以估计若干未知参数。如果这 \(N\)
个测量是相关的，在构造 \(\chi^2\) 时需要用到协方差矩阵的逆
\(V^{-1}\)。一般来说，我们先通过某种办法（比如蒙特卡罗计算）得到相关系数矩阵
\(\rho_{ij}=V_{ij}/(\sigma_i\sigma_j)\)（其中
\(i,j=1,2,\ldots,N\)），然后再计算得到协方差矩阵的逆。

\begin{enumerate}
\item 注意到对于有效估计量，协方差矩阵的逆正比于样本容量
\(n\)。证明：在此条件下，相关系数矩阵与样本容量 \(n\) 无关。

\item 证明协方差矩阵的逆为

\[
(V^{-1})_{ij}=\frac{(\rho^{-1})_{ij}}{\sigma_i\sigma_j}.
\]

【提示：从下面等式出发

\[
\delta_{ij}=\sum_k(V^{-1})_{ik}V_{kj}=\sum_k(V^{-1})_{ik}\rho_{kj}\sigma_k\sigma_j
\]

对式（7.10）两边都乘以
\(\rho^{-1}\)，并对适当的指标求和即可得到式（7.9）。】
\end{enumerate}
\begin{solution}
若协方差矩阵写成

\[
V_{ij}=\rho_{ij}\sigma_i\sigma_j,
\]

其中 \(\rho\) 是相关系数矩阵，则

\begin{enumerate}
\item 若估计量是有效的

通常 \(V\propto 1/n\)。于是 \(\sigma_i\propto n^{-1/2}\)，而

\[
\rho_{ij}=\frac{V_{ij}}{\sigma_i\sigma_j}
\]

中的 \(n\) 因子完全约去，因此 \(\rho\) 与样本量 \(n\) 无关。

\item 证明逆矩阵

记 \(D=\mathrm{diag}(\sigma_1,\dots,\sigma_N)\)，则

\[
V=D\rho D.
\]

因此

\[
V^{-1}=D^{-1}\rho^{-1}D^{-1},
\]

即

\[
(V^{-1})_{ij}=\frac{(\rho^{-1})_{ij}}{\sigma_i\sigma_j}.
\]
\end{enumerate}
\end{solution}
\end{question}



\begin{question}
考虑随机变量 \(x\) 的两个部分重叠的样本，它们的样本容量分别为 \(n\) 和
\(m\)，共有部分的样本容量为 \(c\)。假设已知 \(x\) 的方差
\(V[x]=\sigma^2\)。考虑样本均值

\[
y_1=\frac1n\sum_{i=1}^{n}x_i
\]

和

\[
y_2=\frac1m\sum_{i=1}^{m}x_i.
\]

\begin{enumerate}
\item 证明协方差为

\[
\operatorname{cov}[y_1,y_2]=\frac{c\sigma^2}{nm}.
\]

\item 利用 7.6 节的结果，求 \(y_1\) 和 \(y_2\) 的加权平均和方差。
\end{enumerate}
\begin{solution}
两个样本均值为
\[
y_1=\frac1n\sum_{i=1}^n x_i,
\qquad
y_2=\frac1m\sum_{j=1}^m x_j.
\]
只有共有的 \(c\) 个观测会同时出现在两个求和中，因此
\[
\operatorname{cov}(y_1,y_2)
=\frac{1}{nm}\,c\sigma^2
=\frac{c\sigma^2}{nm}.
\]
并且
\[
V[y_1]=\frac{\sigma^2}{n},
\qquad
V[y_2]=\frac{\sigma^2}{m}.
\]
于是
\[
V=\sigma^2
\begin{pmatrix}
1/n & c/(nm)\\[3pt]
c/(nm)&1/m
\end{pmatrix}.
\]

对同一真均值 \(\mu\) 的两个相关估计量，最佳线性无偏组合为
\[
\hat\mu=\frac{\mathbf 1^T V^{-1}y}{\mathbf 1^T V^{-1}\mathbf 1}.
\]
代入上面的 \(V\) 并化简，得到
\[
\hat\mu
=\frac{(n-c)y_1+(m-c)y_2}{n+m-2c}.
\]
其方差为
\[
V[\hat\mu]
=\frac{1}{\mathbf 1^T V^{-1}\mathbf 1}
=\sigma^2\frac{nm-c^2}{nm(n+m-2c)}.
\]
检查两个极限：若 \(c=0\)，则 \(\hat\mu=(ny_1+my_2)/(n+m)\)，
方差为 \(\sigma^2/(n+m)\)；若第一个样本完全包含在第二个样本中
\((c=n<m)\)，则 \(\hat\mu=y_2\)，方差为 \(\sigma^2/m\)。
\end{solution}
\end{question}



\begin{question}
天文学家托勒密（Claudius
Ptolemy）利用圆盘做过光折射的实验。他把圆盘的一半浸入水中，圆心正好位于水面处，如图
7.2 所示。大约公元 140 年，Ptolemy 对 8 组不同的入射角 \(\theta_i\)
测量了相应的折射角 \(\theta_r\)，结果如表 7.2
所示。本练习中，我们认为入射角已知且误差可以忽略，而把折射角看作标准差为
\(\sigma=\frac12^\circ\)
的高斯随机变量。（这是一个合理的假设，记录的角度精确到最邻近的半度。注意，我们可以将
\(\theta_i\) 的误差吸收到 \(\theta_r\) 的有效误差中。）
\pyfig[0.62\linewidth]{figures/fig_07_02_ptolemy_schematic.png}{图 7.2：Ptolemy 折射实验的角度定义示意图。}


\begin{center}
\begin{tabular}{cc}
\toprule
\(\theta_i\) & \(\theta_r\) \\
\midrule
10 & 8 \\
20 & \(15\frac12\) \\
30 & \(22\frac12\) \\
40 & 29 \\
50 & 35 \\
60 & \(40\frac12\) \\
70 & \(45\frac12\) \\
80 & 50 \\
\bottomrule
\end{tabular}
\par\vspace{0.25em}{\small 表 7.2: Ptolemy 折射实验的入射角和折射角，单位均为度。\par}
\end{center}

\begin{enumerate}
\item 直到 17 世纪才发现正确的折射定律，在此之前，通常的假设是

\[
\theta_r=\alpha\theta_i,
\]

然而 Ptolemy 更喜欢用下面的形式

\[
\theta_r=\alpha\theta_i-\beta\theta_i^2.
\]

对这两种不同的假设，求参数的最小二乘估计量，并计算最小 \(\chi^2\)
值。评论一下两个假设的拟合优度。是否可以相信所有的数据都是从实际测量得来的？

\item 1621 年 Snell 发现了折射定律

\[
\theta_r=\sin^{-1}\left(\frac{\sin\theta_i}{r}\right),
\]

其中 \(r=n_r/n_i\) 为两种介质的折射率之比。求 \(r\)
的最小二乘估计量并计算出最小 \(\chi^2\) 值。评价对 \(\theta_r\) 作
\(\sigma=\frac12^\circ\) 假设的合理性。
\end{enumerate}
\begin{solution}
数据为

\[
(\theta_i,\theta_r)=
(10,8),(20,15.5),(30,22.5),(40,29),(50,35),(60,40.5),(70,45.5),(80,50),
\]

且把 \(\theta_r\) 看作标准差为 \(\sigma=0.5^\circ\) 的高斯变量。

\begin{enumerate}
\item 经验模型比较

\begin{enumerate}
\item
  线性模型 \[
  \theta_r=\alpha\theta_i.
  \]

  得到 \[
  \hat\alpha=0.666176,
  \qquad
  \chi^2_{\min}=134.65,
  \qquad
  n_{\rm dof}=7,
  \qquad
  P\approx6.7\times10^{-26}.
  \]
\item
  Ptolemy 喜欢的经验式 \[
  \theta_r=\alpha\theta_i-\beta\theta_i^2.
  \]

  得到 \[
  \hat\alpha=0.825000,
  \qquad
  \hat\beta=2.5000\times10^{-3},
  \] \[
  \chi^2_{\min}\approx0,
  \qquad
  n_{\rm dof}=6,
  \qquad
  P\approx1.
  \]
\end{enumerate}

这个结果看上去``好得过头''，更像是数据本身被这种简单经验式近似生成或已经经过明显取整；因此不能机械地把
\(P=1\) 理解为``模型绝对正确''。

\item Snell 定律

设

\[
\theta_r=\sin^{-1}\!\left(\frac{\sin\theta_i}{r}\right),
\]

其中 \(r=n_r/n_i\)。数值最小化得到

\[
\hat r=1.31162,
\qquad
\chi^2_{\min}=14.00,
\qquad
n_{\rm dof}=7,
\qquad
P=0.0511.
\]

这已经接近通常的可接受边界，而且 \(r\approx1.31\)
对空气-水界面是很合理的数值。若把 \(\sigma=0.5^\circ\)
看作记录精度给出的``有效标准差''，那么 Snell 定律的拟合并不差。
\pyfig[0.72\linewidth]{figures/fig_07_02_ptolemy_fit.png}{Ptolemy 折射数据及线性、二次经验式、Snell 定律拟合。}\end{enumerate}
\end{solution}
\end{question}



\begin{question}
重新考虑练习 6.9：\(N\) 个独立的泊松变量 \(n=(n_1,\ldots,n_N)\)，均值为
\(\nu=(\nu_1,\ldots,\nu_N)\)，其中均值与某控制变量 \(x\) 有关，

\[
\nu(x)=\theta a(x).
\]

\begin{enumerate}
\item
  首先考虑最小二乘方法（LS），\(\chi^2\) 的分母使用
  \(\sigma_i^2=\nu_i\)。证明 \(\theta\) 的最小二乘估计量为
\end{enumerate}

\[
\hat\theta=\left(\frac{\sum_{i=1}^{N}n_i^2}{\sum_{i=1}^{N}a(x_i)}\right)^{1/2}.
\]

通过对 \(\hat\theta(n)\) 在 \(\nu\)
处进行泰勒展开至第二阶，计算期待值，证明 (7.18) 的偏置为

\[
b=\frac{N-1}{2\sum_{i=1}^{N}a(x_i)}+O(E[(n_i-\nu_i)^3]).
\]

（利用独立泊松变量的协方差
\(\operatorname{cov}[n_i,n_j]=\delta_{ij}\nu_j\)。）

\begin{enumerate}[start=2]
\item
  取 \(\chi^2\) 的分母为
  \(\sigma_i^2=n_i\)，即把观测值作为方差，用修正的最小二乘法（MLS）重复
\begin{enumerate}
\item 中的步骤。证明 \(\theta\) 的最小二乘估计量为
\end{enumerate}

\[
\hat\theta=\frac{\sum_{i=1}^{N}a(x_i)}{\sum_{i=1}^{N}a(x_i)^2/n_i},
\]

并且偏置为

\[
b=-\frac{N-1}{\sum_{i=1}^{N}a(x_i)}+O(E[(n_i-\nu_i)^3]).
\]

将 (a) 和 (b) 得到的偏置与本章第 2 题进行比较。

\begin{enumerate}[start=3]
\item
  利用误差传递，对 LS 和 MLS 两种情况估计 \(\hat\theta\) 的方差。
\end{enumerate}

注意，由于习题 (6.9) 已经证明了 \(\theta\)
的最大似然估计量是无偏的并且方差最小，这里并不推荐最小二乘（LS）和修正的最小二乘（MLS）估计量。然而，对于足够大的数据样本，三个方法是类似的，参见习题
(7.8)。
\end{enumerate}
\begin{solution}
设 \(a_i=a(x_i)\)，\(n_i\sim\operatorname{Pois}(\nu_i)\)，且
\(\nu_i=\theta a_i\)。

\begin{enumerate}
\item LS：\(\sigma_i^2=\nu_i=\theta a_i\)

最小化
\[
\chi^2_{LS}=\sum_i\frac{(n_i-\theta a_i)^2}{\theta a_i}
=\sum_i\left(\frac{n_i^2}{\theta a_i}-2n_i+\theta a_i\right).
\]
对 \(\theta\) 求导并令零，得到
\[
\hat\theta_{LS}
=\left(\frac{\sum_i n_i^2/a_i}{\sum_i a_i}\right)^{1/2}.
\]
按一般的 \(\nu_i=\theta a_i\) 情形，\(n_i^2\) 项必须带有权重 \(1/a_i\)；在所有 \(a_i=1\) 时才化为不含 \(a_i\) 的形式。

在 \(n_i=\nu_i\) 附近二阶展开，利用
\(\operatorname{cov}(n_i,n_j)=\delta_{ij}\nu_i\)，可得
\[
b_{LS}=E[\hat\theta_{LS}]-\theta
=\frac{N-1}{2\sum_i a_i}+O(E[(n_i-\nu_i)^3]).
\]
所以普通 LS 有正偏。

\item MLS：\(\sigma_i^2=n_i\)

最小化
\[
\chi^2_{MLS}=\sum_i\frac{(n_i-\theta a_i)^2}{n_i}
=\sum_i\left(n_i-2\theta a_i+\frac{\theta^2a_i^2}{n_i}\right).
\]
对 \(\theta\) 求导得
\[
\hat\theta_{MLS}
=\frac{\sum_i a_i}{\sum_i a_i^2/n_i}.
\]
同样在 \(n_i=\nu_i\) 附近展开，得到
\[
b_{MLS}
= -\frac{N-1}{\sum_i a_i}+O(E[(n_i-\nu_i)^3]).
\]
因此 MLS 有负偏，且一阶偏置大小约为 LS 的两倍。

\item 方差

对两种估计量在 \(n_i=\nu_i\) 处作一阶误差传播，都有
\[
\left.\frac{\partial\hat\theta}{\partial n_i}\right|_{n=\nu}
=\frac1{\sum_j a_j}.
\]
于是
\[
V[\hat\theta]
\approx\sum_i\left(\frac1{\sum_j a_j}\right)^2V[n_i]
=\frac{\theta}{\sum_i a_i}.
\]
这与第 6.9 题的 MLE 方差相同只是大样本一阶近似；有限样本下 LS 和 MLS 仍分别有上述偏置，因此泊松问题中更推荐直接使用最大似然估计量
\[
\hat\theta_{MLE}=\frac{\sum_i n_i}{\sum_i a_i}.
\]
\end{enumerate}
\end{solution}
\end{question}



\begin{question}
重新考虑 Perrin 关于乳香粒子作为高度的函数（第 6 章第 11 题）。通过最小化

\[
\chi^2(k,\nu_0)=\sum_{i=1}^{N}\frac{(n_i-\nu_i(k,\nu_0))^2}{\sigma_i^2},
\]

求玻尔兹曼常数 \(k\)（或者等价于阿伏伽德罗常数 \(N_A=R/k\)）和系数
\(\nu_0\) 的最小二乘估计量。

\begin{enumerate}
\item
  取 \(n_i\) 的标准差 \(\sigma_i\) 为
  \(\sqrt{\nu_i}\)（通常的最小二乘法）。
\item
  取 \(\sigma_i\) 为 \(\sqrt{n_i}\)（修正的最小二乘法）。
\end{enumerate}

将 (a) 和 (b) 得到的估计量与习题 (6.11) 中最大似然估计量进行比较。
\begin{solution}
沿用 6.11 的 Perrin 数据：

\[
(z, n)= (0,1880),(6,940),(12,530),(18,305),
\]

理论模型为

\[
\nu(z)=\nu_0\exp\!\left[-\frac{4\pi r^3\Delta\rho g z}{3kT}\right].
\]

\begin{enumerate}
\item 取 \(\sigma_i=\sqrt{\nu_i}\)

最小化

\[
\chi^2(k,\nu_0)=\sum_i\frac{(n_i-\nu_i)^2}{\nu_i}
\]

数值结果：

\[
\hat\nu_0=1845.52,
\qquad
\hat k=1.19946\times10^{-16}\,\mathrm{erg/K},
\qquad
\hat N_A=\frac{R}{\hat k}=6.936\times10^{23}\,\mathrm{mol}^{-1}.
\]

并且

\[
\chi^2_{\min}=4.569,
\qquad
n_{\rm dof}=2,
\qquad
P=0.102.
\]

\item 取 \(\sigma_i=\sqrt{n_i}\)

最小化

\[
\chi^2(k,\nu_0)=\sum_i\frac{(n_i-\nu_i)^2}{n_i}
\]

得到

\[
\hat\nu_0=1843.83,
\qquad
\hat k=1.19710\times10^{-16}\,\mathrm{erg/K},
\qquad
\hat N_A=6.950\times10^{23}\,\mathrm{mol}^{-1},
\]

\[
\chi^2_{\min}=4.619,
\qquad
P=0.0993.
\]

\end{enumerate}

与第 6 章第 11 题的 MLE 比较。

MLE 给出

\[
\hat\nu_0=1844.94,
\qquad
\hat k=1.19870\times10^{-16}\,\mathrm{erg/K},
\qquad
\hat N_A=6.941\times10^{23}\,\mathrm{mol}^{-1}.
\]

三种方法在这组数据上的中心值很接近；因为数据量虽少，但每个 bin
的计数较大，泊松分布已接近高斯近似。不过从原则上看，MLE
仍然更自然，因为它直接使用了泊松似然，而两种最小二乘都引入了额外的近似。
\end{solution}
\end{question}

